\chapter{Diffusion models}

\textbf{NB:} This is based on \cite{sohldickstein2015deep}, \cite{ho2020denoising}, \cite{song2019generative} and \cite{song2021scorebased}.


\paragraph{Motivation.}

The flow-matching framework seen in the last chapter can be considered as an implicit generative model, in the sense that the density is not parametrised directly, but rather induced by a sample-and-transport generative procedure. In this chapter we take an alternative approach to implicit generative modelling which, in the same spirit as VAEs and flow-based methods, starts from a source Gaussian distribution and transforms the samples to match a target density that is only specified by an observed dataset. Though we will see that there is a clear connection between the flow-based approach and the diffusion approaches to be reviewed in this chapter, their main difference is conceptual: while flow methods rely on deterministic dynamics defined by an ODE, diffusion models build upon stochastic dynamics, meaning that the path from source to target has an equivalent stochastic differential equation (SDE) formulation. There are two approaches to materialise this rational as we will see in this chapter: score-based models and diffusion models.

\section{Score-based models}

To bypass the challenge of direct parametrisation of the target density, we instead consider the \emph{score function} of the distribution defined as follows.

\begin{definition}[Score function]
\label{def:score_function}
Let $p(\x)$ be a probability density on $\R^D$. The \emph{score function} of the distribution is defined as the gradient of its log-density
\begin{equation}
    \label{eq:score_function}
    s(\x) = \nabla_\x \log p(\x).
\end{equation}
\end{definition}

\begin{remark}[Intuition]
The vector $s(\x)\in\R^D$ points in the direction of the steepest increase of the log-density, that is, toward nearby regions where the density becomes larger. Furthermore, its magnitude reflects how rapidly the density changes around $\x$. Intuitively, if we imagine particles distributed according to $p(\x)$, the score indicates how a particle located at $\x$ should move in order to reach regions of higher probability density.
\end{remark}

Targetting the score, rather than the density, has clear advantages as discussed in the following remark. 

\begin{remark}[Unrestrictedness of the score]
    By modelling the score, we lift the restriction to model a normalised function. Indeed, one can write 
    \[
    p(\x) = \frac{1}{z}\tilde{p}(\x), \qquad z = \int \tilde{p}(\x)\d\x,
    \]
    where $\tilde{p}(\x)$ is the unnormalised version of the density $p(\x)$ which is, arguably, easier to model. This is because we can use an unrestricted class of models, such as neural networks, which would be unfeasible to model densities, because the normalisation restriction could not be directly enforced. The score of such a distribution is
    \[
    s(\x) = \nabla_\x\log \tilde{p}(\x) - \cancelto{=0}{\nabla_\x z},
    \]
    meaning that the score does not depend on the normalising constant. 
\end{remark}

Once the score is learnt, via, e.g., a neural network with an appropriate training procedure to be discussed later, we need to implement a methodology to produce samples $\x\sim p(\x)$ by only having access to the score in eq.~\eqref{eq:score_function}; and not to $p(\cdot)$. We will consider the following method. 


\begin{definition}[Langevin sampler]
\label{def:langevin_sampler}
Let $p(\x)$ be a target density with score function $s(\x)=\nabla_\x \log p(\x)$.
The Langevin sampler generates samples from $p(\x)$ by iteratively updating
\begin{equation}
    \x_{k+1} = \x_k + \frac{\epsilon}{2}s(\x_k) + \sqrt{\epsilon}\,\eta_k,
\end{equation}
where $\epsilon>0$ is a step size and $\eta_k \sim \cN(0,\I)$ is Gaussian noise.
\end{definition}

The update in Def.~\ref{def:langevin_sampler} can be interpreted as a random walk with a correction in the direction of the score, or equivalently, a stochastic gradient ascent step on the log-density combined with Gaussian noise injection. This can be understood as a balance between the score term, which pushes particles toward regions of higher probability mass, and  the noise term, which ensures exploration of the state space.

\begin{remark}[Stationary distribution of Langevin dynamics]
Under appropriate conditions on the step size and regularity of the score function, the Markov chain defined in Def.~\ref{def:langevin_sampler} admits $p(\x)$ as its stationary distribution. Therefore, as $k \to \infty$, the distribution of
$\x_k$ converges to $p(\x)$. In practice, after a sufficiently large number of iterations, the samples $\x_k$ can be regarded as approximate draws from the target distribution.
\end{remark}

\paragraph{Difficulties in score-based generative modelling.} Two main challenges arise in this setting. The first one stems from the so-called \emph{data manifold hypothesis}, which states that, in practice, the data lives in a low-dimensional manifold within the ambient space. As a consequence, even if within the data manifold the density behaves well, as soon as we leave the manifold the probability goes sharply to zero, meaning that the measure is not absolutely continuous wrt the Lebesgue measure and thus the density does not exist. As a consequence, the score is not defined as it involves a gradient taken in the ambient space. Therefore, we cannot attempt to learn the score from data.

A second challenge stems from the data scarcity: to sample using Langevin dynamics, it is required that the score is appropriately estimated on the entire ambient space, so that the sampler can navigate the space from areas of low probability to areas of high probability. However, in the areas of low probability, there will be insufficient samples to train a score function reliably, and this the sampling procedure will not be accurate and the chain will not mix well.  

\begin{remark}[Key idea of the approach]
The two challenges described above will be addressed by introducing a smoothed version of the data distribution. By perturbing the data with Gaussian noise, the resulting distribution becomes well-defined over the entire ambient space, ensuring that the score exists everywhere. This perturbation will also allow us to construct a synthetic dataset from which the score can be learned.
\end{remark}


\paragraph{Denoising score matching.} In a way that is similar to the flow-matching objective, a parametrised (e.g., neural network) model of the score could be trained by minimising the following so called score-matching loss
\begin{equation}
    \label{eq:SM_loss1}
    \cL_{\text{SM}}(\theta) = 
    \frac{1}{2}
    \mathbb{E}_{p_\text{data}(\x)}||s_\theta(\x) - \nabla_\x\log p_\text{data}(\x)||_2^2.
\end{equation}
Note that this loss is intractable, as it requires evaluating of the unknown data distribution. This can be bypassed considering the equivalent expression as noted by \cite[Thm.~1]{hyvarinen2005score}
\begin{equation}
     \label{eq:SM_loss2}
    \cL_{\text{SM}}(\theta) 
    =
    \mathbb{E}_{p_\text{data}(\x)}\left[\tr{(\nabla_\x s_\theta(\x))} + \frac{1}{2}||s_\theta(\x)||_2^2\right] + \text{constant},
\end{equation}
where the constant does not depend on the model parameter $\theta$. The formulation in eq.~\eqref{eq:SM_loss2} is now tractable since it only depends on the unknown data density via samples, meaning that it can be approximated via Monte Carlo. However, trace of the Jacobian matrix $\nabla_{\x}s_\theta(\x)\in\R^{D\times D}$ has a computational complexity scaling as $\mathcal O(D^2)$, which makes it prohibitively expensive for high-dimensional data such as images. 

An alternative objective can be obtained by first smoothing the data distribution with a corruption kernel \(q_\sigma(\tilde{\x}\mid \x)\), and then applying score matching to the resulting perturbed density
\begin{equation}
    \label{eq:perturbed_marginal}
    q_\sigma(\tilde{\x})=\int q_\sigma(\tilde{\x}\mid \x)\,p_{\mathrm{data}}(\x)\,\d\x.    
\end{equation}
As noted by \cite{vincent2011connection}, applying the score-matching objective in eq.~\eqref{eq:SM_loss1} to the perturbed data distribution is equivalent, up to an additive constant independent of \(\theta\), to the denoising score matching objective and it is given by
\begin{equation}
    \label{eq:DSM_loss}
    \cL_{\text{DSM}}(\theta) =
    \frac{1}{2}
    \mathbb{E}_{p_{\text{data}}(\x)\, q_\sigma(\tilde{\x}\mid\x)}
    \left\|
    s_\theta(\tilde{\x}) -
    \nabla_{\tilde{\x}} \log q_\sigma(\tilde{\x}\mid\x)
    \right\|_2^2,
\end{equation}
meaning that the minimiser of eq.~\eqref{eq:DSM_loss} satisfies $s_\theta(\tilde\x)^*=\nabla_{\tilde\x}\log q_\sigma(\tilde\x)$. 

\begin{remark}[Validity of the denoising score matching objective]
If the perturbation kernel $q_\sigma(\tilde{\x}\mid\x)$ is sufficiently concentrated around $\x$ (e.g.\ small $\sigma$ in the Gaussian case), then the perturbed marginal
$q_\sigma(\tilde{\x})$ in eq.~\eqref{eq:perturbed_marginal} satisfies
\[
q_\sigma(\tilde{\x}) \approx p_{\mathrm{data}}(\tilde{\x}).
\]
Consequently, the score of the perturbed density approximates the score of the data distribution, and the minimiser of $\cL_{\mathrm{DSM}}(\theta)$ can be regarded as a reliable approximation of the minimiser of the original score-matching objective $\cL_{\mathrm{SM}}(\theta)$.
\end{remark}

Using the perturbation kernel $q_\sigma(\tilde{\x}\mid\x)$, the DSM objective above is theoretically appealing in the sense that the density and the score are well defined on the entire ambient space due to the smoothing operation. Intuitively, using a concentrated perturbation allows for a reliable approximation of the data distribution, while a wide support kernel would populate the low-probability regions in the ambient space, thus facilitating the implementation of the Langevin sampler. 

\paragraph{Noise-conditional score network.} 

Understanding the perturbation kernel as the injection of (Gaussian) noise with variance $\sigma^2$, \cite{song2021scorebased} builds on the observations above to propose a noise-indexed neural network to estimate the score for varying levels of noise defined as follows. 

Let $\{\sigma_i\}_{i=1}^L$ be a geometric sequence satisfying $\sigma_1 > \sigma_2 > \cdots > \sigma_L>0$, and choose a Gaussian perturbation kernel to construct $L$ noise-injected versions of the data distribution, with the $i-$th level perturbed density given by 
\begin{equation}
    q_{\sigma_i} (\x)= \int p_{\mathrm{data}}(\x') \cN(\x \mid \x',\sigma_i^2 \I_D).
\end{equation}
Furthermore, the sequence $\{\sigma_i\}_{i=1}^L$ is chosen so that $\sigma_1$ is large enough to mitigate all the issues discussed above and $\sigma_L$ is small enough so that $q_{\sigma_L} (\x) \approx p_{\mathrm{data}}(\x)$.

Since the score of the Gaussian perturbation kernel is linear, for each noise level $\sigma$ we consider the following DSM target based on eq.~\eqref{eq:DSM_loss} given by 
\begin{equation}
    \label{eq:DSM_loss_noise_level}
    l(\theta;\sigma) = 
    \frac12 \mathbb{E}_{\x\sim p_{\mathrm{data}}(\x), \tilde\x\sim \cN(\x,\sigma^2\I_D)}
    \left(\left\|
        s_\theta(\tilde{\x},\sigma) + \frac{\tilde{\x} - \x}{\sigma^2}
    \right\|^2_2\right),
\end{equation}
and then aggregate them across all noise levels
\begin{equation}
    \cL_{\text{NCSN}}(\theta, \{\sigma_i\}_{i=1}^L) = \frac1L \sum_{i=1}^{L} \lambda(\sigma_i) l(\theta;\sigma_i),
\end{equation}
where $\lambda(\sigma) = \sigma^2$  ensures that $ \lambda(\sigma) l(\theta;\sigma)$ does not depend on $\sigma$.

\begin{remark}[Preparing data for training]
    Training the noise-conditional score network (NCSN) based on the above loss, does not require to construct and store $L$ separate datasets explicitly. Instead, at each training iteration one samples a datapoint $\x\sim p_{\mathrm{data}}$, then samples a noise level $\sigma_i$ uniformly from the sequence $\{\sigma_i\}_{i=1}^L$, to generate a perturbed sample through
    \[
    \tilde{\x} = \x + \sigma_i \z, \qquad \z \sim \cN(0,\I_D).
    \]
    The pair $(\tilde{\x},\sigma_i)$ is then used to evaluate the objective in eq.~\eqref{eq:DSM_loss_noise_level}, with the perturbed training examples generated during training to construct a Monte Carlo estimate of the loss, This way, there is no need to precompute and store noisy copies of the dataset.
\end{remark}



Lastly, after the NCSN is trained, \cite{song2021scorebased} proposes an \emph{annealed} variant of Langevin dynamics, inspired by simulated annealing. The procedure begins by drawing an initial sample from a Gaussian distribution, which coincides with the highest-noise perturbed density,
\[
\x^{(0)} \sim \mathcal N(0,\sigma_1^2 I_D).
\]
Starting from this initial point, Langevin dynamics is run for several iterations targeting the distribution \(q_{\sigma_1}\). The resulting sample is then used to initialise Langevin dynamics targeting the next distribution \(q_{\sigma_2}\), corresponding to a smaller noise level. This procedure is repeated sequentially along the noise schedule \(\sigma_1 > \sigma_2 > \cdots > \sigma_L\), each time using the last sample from the previous stage as the initial condition for the next one. The final stage targets \(q_{\sigma_L}\), which corresponds to the smallest noise level and provides an approximation of the data distribution. By gradually decreasing the noise level, the sampler goes from a coarse exploration of the ambient space toward a reliable approximation of the data distribution.


\begin{algorithm}[H]
\caption{Training a noise-conditional score network (NCSN)}
\label{alg:ncsn_training}
\begin{algorithmic}[1]
\Require Dataset \(\mathcal D=\{\x^{(n)}\}_{n=1}^N\), noise levels \(\{\sigma_i\}_{i=1}^L\), weights \(\lambda(\sigma_i)\), learning rate \(\eta\)
\Ensure Trained parameters \(\theta\)

\State Initialise network parameters \(\theta\)
\While{not converged}
    \State Sample a minibatch \(\{\x_b\}_{b=1}^B\) from \(\mathcal D\)
    \For{$b=1,\ldots,B$}
        \State Sample \(i_b \sim \mathrm{Unif}\{1,\ldots,L\}\)
        \State Draw \(\z_b \sim \mathcal N(0,I_D)\)
        \State Set \(\tilde{\x}_b = \x_b + \sigma_{i_b}\z_b\)
        \State Set \(\t_b = -(\tilde{\x}_b-\x_b)/\sigma_{i_b}^2\)
    \EndFor
    \State Compute
    \[
    \widehat{\mathcal L}(\theta)
    =
    \frac{1}{2B}
    \sum_{b=1}^B
    \lambda(\sigma_{i_b})
    \left\|
    s_\theta(\tilde{\x}_b,\sigma_{i_b})-\t_b
    \right\|_2^2
    \]
    \State Update \(\theta \leftarrow \theta - \eta \nabla_\theta \widehat{\mathcal L}(\theta)\)
\EndWhile
\State \Return \(\theta\)

\end{algorithmic}
\end{algorithm}

\begin{algorithm}[H]
\caption{Annealed Langevin dynamics sampling with a trained NCSN}
\label{alg:ald_sampling}
\begin{algorithmic}[1]
\Require Trained score network \(s_\theta(\x,\sigma)\), noise levels \(\sigma_1>\sigma_2>\cdots>\sigma_L\), step sizes \(\{\alpha_i\}_{i=1}^L\), numbers of steps \(\{T_i\}_{i=1}^L\)
\Ensure Approximate sample \(\x\) from the data distribution

\State Draw initial sample \(\x \sim \mathcal N(0,\sigma_1^2 I_D)\)
\For{$i=1,\ldots,L$}
    \For{$t=1,\ldots,T_i$}
        \State Draw \(\z \sim \mathcal N(0,I_D)\)
        \State Update
        \[
        \x \leftarrow \x + \frac{\alpha_i}{2}s_\theta(\x,\sigma_i) + \sqrt{\alpha_i}\,\z
        \]
    \EndFor
\EndFor
\State \Return \(\x\)

\end{algorithmic}
\end{algorithm}



\section{Denoising diffusion probabilistic models}



Let us consider a variable of interest $\x_0\in\R^D$, such that $\x_0\sim q(\x_0)$, with $q(\cdot)$ its true unknown probability density. Starting from $\x_0$, we can construct a sequence of random variables $\x_1,\ldots,\x_T\in\R^D$, according to the following \emph{forward process} given by Gaussian perturbation kernels: 
\begin{align}
    q (\x_{1:T} \mid \x_0) &= \prod_{t=1}^T q(\x_t \mid \x_{t-1})
    \label{eq:diff_q_factor}\\
    q(\x_t \mid \x_{t-1}) &= \cN(\x_t; \sqrt{1-\beta_t}\x_{t-1},\beta\I), \label{eq:diff_q_def}
\end{align}
where the parameters $\{\beta_t\}_{t=1}^T$ are fixed and chosen so that, given $\x_0$, the sequence $\x_1,\ldots,\x_T$ becomes progressively distributed according to a standard Gaussian, that is, the terminal variable $\x_T$ is expected to be distributed close to $\cN(0,\I)$. 

To see this, note that for an arbitrary $t\in\{1,\ldots,T\}$, the variable $\x_t$ can be directly sampled conditional to $\x_0$ according to 
\begin{equation}
    q(\x_t \mid \x_0) = \cN(\x_t; \sqrt{\bar\alpha_t}\x_0, (1-\bar\alpha_t)\I),   
    \label{eq:noisy_sample_from_x0} 
\end{equation}
where $\alpha_t = 1-\beta_t$, and $\bar\alpha_t = \prod_{s=1}^{t}\alpha_s$. Therefore, for a sufficiently long sequence of steps with small $\beta_t$'s, we have $\bar\alpha_T\approx 0$ and $q(\x_t \mid \x_0) \approx \cN(0,\I)$.

\begin{remark}[True model]
    Observe that the model $q(\x_{0:T}) = q(\x_{1:T} \mid \x_0)q(\x_{0})$ is the true distribution of the process $\x_{0:T}$, and that, if a set of observations $\x_0$ were available, we could sample from the rest of the sequence $\x_{1:T}$. However, this formulation does not allow for sampling from the data distribution $q(\x_{0})$, which is unknown.  
\end{remark}

Given that our ultimate goal is to sample from the data distribution, we consider an approximating  parametric model for the sequence $\x_{0:T}$, given by $p_\theta(\x_{0:T})$. Since we know the distribution of the terminal variable $\x_T$ is approximately a standard Gaussian, we will choose $p_\theta(\x_T) = p(\x_T)= \cN(\x_T; 0, \I)$ and then define the rest of the approximating model as a \emph{reverse process} given by 
\begin{align}
    p_\theta (\x_{0:T}) &= p(\x_T) \prod_{t=1}^T p_\theta(\x_{t-1} | \x_t),\label{eq:diff_p_factor}\\
    p_\theta(\x_{t-1} | \x_t) &= \cN(\x_{t-1}; \mu_\theta(\x_t,t), \Sigma_\theta(\x_t,t)).
    \label{eq:diff_p_def}
\end{align}

Observe that $p_\theta(\x_{0:T})$ defines a latent-variable model, where the approximation to the data distribution is given by 
\[
p_\theta (\x_0) = \int p_\theta(\x_{0:T})\d\x_{1:T},
\]
and, as opposed to previous latent-variable model we have seen, both ambient and latent  variables here are of the same dimension, i.e.,  $\x_{0:T}\in\R^D$.

\begin{remark}[Approximating model]
    \label{rem:approximating_model}
    Unlike $q(\x_0)$, sampling from $p_\theta(\x_0)$ using the reverse process in eqs.~\eqref{eq:diff_p_factor}-\eqref{eq:diff_p_def} is straightforward since the initial distribution $\x_T\sim p(\x_T) = \cN(0,\I)$ is Gaussian. Therefore, we will train $p_\theta(\x_{0:T})$ so that $p_\theta(\x_{0})$ matches $q(\x_{0})$
\end{remark}


\begin{remark}[Hierarchical VAE analogy]
    Observe that the model introduced above has a \emph{deep} (or hierarchical) VAE structure. 
    To see this, consider first the case $T=1$. The generative model becomes
    \[
        p_\theta(\x_0,\x_1) = p(\x_1)p_\theta(\x_0 \mid \x_1),
    \]
    where $p(\x_1)=\cN(0,\I)$ and $p_\theta(\x_0 \mid \x_1)$ is Gaussian with parameters
    given by the networks $\mu_\theta$ and $\Sigma_\theta$. The variational distribution is given by the forward
    diffusion kernel
    \[
        q(\x_1 \mid \x_0) = \cN(\x_1;\sqrt{1-\beta_1}\x_0,\beta_1\I),
    \]
    which plays the role of the encoder in a VAE. Notice, however, that in this case the encoder $q(\x_1|\x_0)$ is fixed rather than learned.

    For general $T$, a hierarchical structure arises: the model can be interpreted as a sequence of VAEs where each layer progressively corrupts (or denoises) the representation produced by the previous one. See \cite{luo2022understanding} for an in-depth presentation of hierarchical VAEs and diffusion models. See a diagram in Fig.~\ref{fig:HVAE}. 
\end{remark}

\begin{figure}[t]
  \centering
  \includegraphics[width=0.9\textwidth]{img/week9_HVAE.JPG}
  \caption{Diagram of the presented hierarchical VAE analogy. Taken from \protect\cite{luo2022understanding}}
  \label{fig:HVAE}
\end{figure}


\begin{remark}[Model expressivity]
The encoder distributions $q(\x_t|\x_{t-1})$ are fixed Gaussian kernels. In contrast to standard VAEs, this does not significantly limit the model expressivity, since the role of the forward process is simply to transform the data distribution into a Gaussian, and the construction above ensures that. The flexibility required to model the data distribution, therefore, stems from the reverse process $p_\theta(\x_{t-1}|\x_t)$, which must learn how to invert this \emph{Gaussianisation} recursively.
\end{remark}

\paragraph{Variational objective.}

In Remark~\ref{rem:approximating_model}, it was stated that the model $p_\theta(\x_{0:T})$ should be trained so that its marginal $p_\theta(\x_0)$ approximates the data distribution $q(\x_0)$. For a fixed observed datapoint $\x_0$, the variables $\x_{1:T}$ are latent, and the forward diffusion process in eq.~\eqref{eq:diff_q_factor}
provides a natural variational distribution. Therefore, applying
Jensen's inequality yields the evidence lower bound (ELBO)
\begin{align}
\elbo(\x_0)
&=
\mathbb{E}_{q(\x_{1:T}\mid \x_0)}
\left(
\log \frac{p_\theta(\x_{0:T})}{q(\x_{1:T}\mid \x_0)}
\right)
\le \log p_\theta(\x_0).
\end{align}

Furthermore, incorporating the forward/backward factorisations in eqs.~\eqref{eq:diff_p_factor} and \eqref{eq:diff_q_factor}, we obtain
\begin{align}
    \elbo(\x_0)
    &= \mathbb{E}_{q(\x_{1:T} \mid \x_0)} \left(\log p(\x_T) + \sum_{t=2}^{T}\log\frac{p_\theta (\x_{t-1} \mid \x_t)}{ q(\x_t \mid \x_{t-1})} + \log\frac{p_\theta (\x_{0} \mid \x_1)}{ q(\x_1 \mid \x_{0})} \right).\label{eq:ELBO_DDPM1}
\end{align}


\begin{remark}[ELBO computability vs.\ interpretability]
Though the above expression is tractable in the sense that all the terms are Gaussian and the expectation can be estimated via Monte Carlo, it does not yet reveal how each transition kernel $p_\theta(\x_{t-1}\mid\x_t)$ should match the reverse dynamics of the forward process.

Ideally, we would like to identify terms of the form
\[
\KL{q(\x_{t-1}\mid\x_t)}{p_\theta(\x_{t-1}\mid\x_t)},
\]
which would explicitly compare the learned reverse kernel with the true
reverse conditional of the diffusion process. However, this conditional
depends on the unknown data distribution through
\begin{align*}
q(\x_{t-1}\mid\x_t)
&=
\int q(\x_{t-1}\mid\x_t,\x_0)\,q(\x_0\mid\x_t)\,d\x_0.\\
&= \int q(\x_{t-1}\mid\x_t,\x_0)\,q(\x_{t} \mid \x_0) \frac{q(\x_0)}{q(\x_{t})} \d\x_0.
\end{align*}

\end{remark}

We can introduce a simple reformulation of the terms in eq.~\eqref{eq:ELBO_DDPM1} to make it more interpretable in the sense of the previous remark. By Bayes' rule, we have
\[
q(\x_{t-1}\mid \x_t,\x_0)
=
\frac{q(\x_t\mid \x_{t-1},\x_0)\,q(\x_{t-1}\mid \x_0)}{q(\x_t\mid \x_0)},
\]
however, since $\x_{1:T}$ is Markovian given $\x_0$ under $q(\cdot)$, we have $q(\x_t\mid \x_{t-1},\x_0)=q(\x_t\mid \x_{t-1})$, and thus
\begin{equation}
    q(\x_t\mid \x_{t-1})
=
\frac{q(\x_{t-1}\mid \x_t,\x_0)\,q(\x_t\mid \x_0)}{q(\x_{t-1}\mid \x_0)}.
\label{eq:bayes_in_DDPM}
\end{equation}

\begin{remark}[Closed form of eq.~\eqref{eq:bayes_in_DDPM}]
    Conditional on $\x_0$, the variational denoiser is Gaussian and given by 
    \begin{align}
q(\x_{t-1}\mid \x_t, \x_0)
&=
\mathcal{N}\bigl(\x_{t-1}; \tilde{\mu}_t(\x_t,\x_0), \tilde{\beta}_t I \bigr), \label{eq:conditional_denoiser} \\
\tilde{\mu}_t(\x_t,\x_0)
&:=
\frac{\sqrt{\bar{\alpha}_{t-1}}\beta_t}{1-\bar{\alpha}_t}\x_0
+
\frac{\sqrt{\alpha_t}(1-\bar{\alpha}_{t-1})}{1-\bar{\alpha}_t}\x_t, \label{eq:conditional_denoiser_mean}\\
\tilde{\beta}_t
&:=
\frac{1-\bar{\alpha}_{t-1}}{1-\bar{\alpha}_t}\beta_t. \label{eq:conditional_denoiser_var}
\end{align}
We will need this expression later. 
\end{remark}

Incorporating eq.~\eqref{eq:bayes_in_DDPM} into eq.~\eqref{eq:ELBO_DDPM1}, we have 
\begin{align}
    \elbo(\x_0)
     &= \mathbb{E}_{q(\x_{1:T} \mid \x_0)} \left(\log p(\x_T) + \sum_{t=2}^{T}\log\frac{p_\theta (\x_{t-1} \mid \x_t)}{ q(\x_{t-1}\mid \x_t,\x_0)}\frac{q(\x_{t-1}\mid \x_0)}{q(\x_t\mid \x_0)} + \log\frac{p_\theta (\x_{0} \mid \x_1)}{ q(\x_1 \mid \x_{0})} \right)\nonumber\\
     &= \mathbb{E}_{q(\x_{1:T} \mid \x_0)} \left(\log p(\x_T) 
     + \sum_{t=2}^{T}\left(\log\frac{p_\theta (\x_{t-1} \mid \x_t)}{ q(\x_{t-1}\mid \x_t,\x_0)}
     + \log\frac{q(\x_{t-1}\mid \x_0)}{q(\x_t\mid \x_0)} \right)
     + \log\frac{p_\theta (\x_{0} \mid \x_1)}{ q(\x_1 \mid \x_{0})} \right)\nonumber\\
     &= \mathbb{E}_{q(\x_{1:T} \mid \x_0)} \left(
    \log\frac{p(\x_T)}{q(\x_T\mid \x_0)}  
     + \sum_{t=2}^{T}\log\frac{p_\theta (\x_{t-1} \mid \x_t)}{ q(\x_{t-1}\mid \x_t,\x_0)}
     + \log p_\theta (\x_{0} \mid \x_1) \label{eq:ELBO_DDPM2}
     \right),
\end{align}
where in the last expression we used the fact that the terms $q(\x_t\mid \x_0)$ cancel each other due to the telescopic sum. 

Recall that the objective in eq.~\eqref{eq:ELBO_DDPM2} corresponds to a single datapoint $\x_0$. To obtain the ELBO corresponding to the full data distribution we take the expectation with
respect to $q(\x_0)$:
\[
\elbo = \mathbb{E}_{q(\x_0)}[\elbo(\x_0)].
\]
Identifying the joint distribution of the forward diffusion process
\(
q(\x_{0:T}) = q(\x_0)q(\x_{1:T}\mid \x_0),
\)
the objective becomes
\begin{align}
\elbo
&=
\mathbb{E}_{q(\x_{0:T})}
\left(
\log\frac{p(\x_T)}{q(\x_T\mid \x_0)}
+
\sum_{t=2}^{T}\log\frac{p_\theta (\x_{t-1} \mid \x_t)}{ q(\x_{t-1}\mid \x_t,\x_0)}
+
\log p_\theta (\x_{0} \mid \x_1)
\right).
\label{eq:ELBO_DDPM3}
\end{align}

Observe that the first two terms can be expressed as KL divergences by extracting the appropriate marginal  from the expectation. This is indeed possible since we can write 
\[
q(\x_{0:T})
=
q(\x_0)\,q(\x_t\mid \x_0)\,q(\x_{t-1}\mid \x_t,\x_0)\,
q(\x_{1:t-2},\x_{t+1:T}\mid \x_{t-1},\x_t,\x_0).
\]
Therefore, we obtain
\[
\elbo = \mathbb{E}_{q(\x_{0:T})} \left[\sum_{t=0}^{T} L_t\right]
\]
where
\begin{align*}
L_T
&=
-
\KL{q(\x_T\mid \x_0)}{p(\x_T)},\\[4pt]
L_{t-1}
&=
-
\KL{q(\x_{t-1}\mid \x_t,\x_0)}{p_\theta(\x_{t-1}\mid \x_t)},
\qquad t=2,\dots,T,\\[4pt]
L_0
&=
\log p_\theta(\x_0\mid \x_1).
\end{align*}

\begin{remark}[Interpretation of the terms in the ELBO]
    The term $L_T$ matches the terminal forward distribution to the prior, the terms $L_{t-1}$ for $t=2,\dots,T$ match each learned reverse kernel to the corresponding reverse conditional of the forward process, and $L_0$ corresponds to a reconstruction likelihood of the data from the first latent variable $\x_1$.
\end{remark}

\paragraph{Design of the denoising autoencoder.} Based on the loss above, we now establish particular choices for the design of the reverse generative process $p_\theta$. 

First, note that the $L_T$ term measures the discrepancy between $p(\x_T)$, which is chosen as a standard Gaussian, and $q(\x_T \mid \x_0)$ which is \emph{close to} a standard Gaussian based on the choice of the schedule $\{\beta_t\}_{t=1}^T$ which is fixed. Therefore, this term has no trainable parameters $\theta$ and thus will be ignored from the optimisation. 

Second, the term $L_0$ measures the reconstruction of the data from the first latent variable and therefore plays the role of a \emph{link} function ensuring that the output of the reverse process matches the format of the data. In \cite{ho2020ddpm}, which focuses on image generation, data are assumed to take integer values in $\{0,\ldots,255\}$ linearly scaled to $[-1,1]$. Therefore, a discrete log-likelihood is obtained in the final step of the reverse process by implementing a discrete decoder derived from a Gaussian distribution, that is, by integrating this Gaussian over the quantization
interval corresponding to each pixel value. As a consequence, and unlike the term $L_T$, this component does involve trainable parameters and therefore contributes to the optimisation.

Third, for the $L_{t-1}$ terms, the backward Gaussian kernel in eq.~\eqref{eq:diff_p_def} can be chosen as 
\[
p_\theta(\x_{t-1} \mid \x_t) = \cN(\x_{t-1} ; \mu_\theta(\x_t,t), \sigma_t^2\I),
\]
with, e.g., $\sigma_t^2 = \beta_t$. With this choice, we have 
\begin{equation}
    L_{t-1} \propto_\theta \mathbb{E}_q(\x_{0:T})\left[ \frac{1}{2\sigma_t^2 }\| \tilde\mu_t(\x_t,\x_0) - \mu_\theta(\x_t,t) \|^2\right],
    \label{eq:L_t-1_with_mus}
\end{equation}
where recall that the symbol $\propto_\theta$ denotes proportionality up to terms independent of $\theta$.

\begin{remark}[Design of $\mu_\theta(\x_t,t)$]
  The denoising model $\mu_\theta(\x_t,t)$ is trained to match the variational denoiser $\tilde\mu_t(\x_t,\x_0)$ in eq.~\eqref{eq:conditional_denoiser_mean}; this provides the necessary insight to design  $\mu_\theta(\x_t,t)$.   
\end{remark}
Recall that  $\tilde\mu_t(\x_t,\x_0)$ is known in close from in eq.~\eqref{eq:conditional_denoiser_mean}, and that $\x_t$ can be expressed according to eq.~\eqref{eq:noisy_sample_from_x0} as $\x_t = \sqrt{\bar\alpha_t}\x_0 + \epsilon\sqrt{1-\bar\alpha_t}$, for $\epsilon\sim\cN(0,\I)$. Therefore, we can write 

\begin{align}
    \tilde{\mu}_t(\x_t,\x_0)
    &=
    \frac{\sqrt{\bar{\alpha}_{t-1}}\beta_t}{1-\bar{\alpha}_t}\x_0
    +
    \frac{\sqrt{\alpha_t}(1-\bar{\alpha}_{t-1})}{1-\bar{\alpha}_t}\x_t\nonumber\\
    &= \frac1{\sqrt{\alpha_t}}\left(\x_t - \frac{\beta_t}{\sqrt{1-\bar\alpha_t}}\epsilon\right),  \label{eq:new_mu_tilde}
\end{align}
where we have used the identities $\beta_t = 1-\alpha_t$ and $\bar\alpha_t = \prod_{s=1}^{t}\alpha_s$.



Since $\mu_\theta(\x_t,t)$ estimates the above expression as a function of $\x_t$, we use the fact that the schedule $\beta_t$ is fixed to consider the following reparametrisation: 

\begin{equation}
    \label{eq:new_mu_theta}
    \mu_\theta(\x_t,t) = \frac1{\sqrt{\alpha_t}}\left(\x_t - \frac{\beta_t}{\sqrt{1-\bar\alpha_t}}\epsilon_\theta(\x_t,t)\right),
\end{equation}
where now $\epsilon_\theta:\R^D\times\R \to \R^D$ is a networks that estimates the noise $\epsilon$ using $\x_t$ and $t$.

Lastly, updating eqs.~\eqref{eq:new_mu_theta} and \eqref{eq:new_mu_tilde} into eq.~\eqref{eq:L_t-1_with_mus}, and expressing again $\x_t$ as a function of the datapoint $\x_0$ so that the expectation is only take wrt $\x_0$ and $\epsilon$, we have the new simplified $L_{t-1}$ loss

\begin{equation}
    L_{t-1} \propto_\theta \mathbb{E}_{\x_0,\epsilon}\left[ \frac{\beta_t^2}{2\sigma_t^2\alpha_t (1-\bar\alpha_t) }\| \epsilon - \epsilon_\theta(\sqrt{\bar\alpha_t}\x_0 + \epsilon\sqrt{1-\bar\alpha_t},t) \|^2\right],
    \label{eq:L_t-1_with_epsilons}
\end{equation}

\paragraph{Denoising diffusion implicit model (DDIM).} A variant of DDPM is DDIM \cite{song2021denoising}, which builds on the observation that DDPM's loss in eq.~\eqref{eq:L_t-1_with_epsilons} is not exclusive to the chosen forward $q(\x_{0:T})$, and thus one can construct a family of non-Markovian noising processes inducing the same training objective, but with different transition densities. Since DDPM samples according to 
\[
\x_{t-1} \sim p_\theta(\x_{t-1}\mid \x_t)
= \cN\!\bigl(\x_{t-1};\mu_\theta(\x_t,t),\sigma_t^2 I\bigr),
\]
it injects Gaussian noise at each iteration. The role of the noise-prediction network $\epsilon_\theta(\x_t,t)$ is therefore not to directly produce the next sample, but rather to predict the noise component in $\x_t$, from which the mean of the reverse Gaussian transition is computed. Consequently, even with a perfect predictor $\epsilon_\theta$, the reverse step in DDPM remains stochastic due to the sampling from the Gaussian distribution.

Conversely, DDIM recycles the trained noise predictor $\epsilon_\theta(\x_t,t)$ (as the training loss is the same) but defines a reverse update given by
\[
\x_{t-1}
=
\sqrt{\bar\alpha_{t-1}}\,\hat{\x}_0(\x_t,t)
+
\sqrt{1-\bar\alpha_{t-1}}\,\epsilon_\theta(\x_t,t),
\qquad
\hat{\x}_0(\x_t,t)
=
\frac{1}{\sqrt{\bar\alpha_t}}
\left(\x_t-\sqrt{1-\bar\alpha_t}\,\epsilon_\theta(\x_t,t)\right),
\]
which corresponds to the deterministic case of the more general DDIM family. Note that the above equation can also be motivated by replacing $\x_0$ (instead of $\x_t$) in eq.~\eqref{eq:new_mu_tilde}. 

Therefore, unlike DDPM, DDIM does not implement a stochastic reverse diffusion during inference, but rather a deterministic procedure. In other words, while DDPM learns a sequence of reverse transition distributions and samples from them, DDIM uses the same trained denoiser to construct a deterministic trajectory from $\x_T$ to $\x_0$.

Critically, the computational gain of DDIM does not come from avoiding the sampling of Gaussian noise, which is negligible compared to evaluating $\epsilon_\theta$. Instead, DDIM uses a subset of timesteps during inference, effectively jumping between noise levels, instead of DDPM which goes through the entire schedule $T,T-1,\dots,1$ (usually $T\approx1000$). Since each step requires evaluating $\epsilon_\theta$, reducing the number of timesteps accelerates sampling while still producing high-quality samples.

\section{Equivalences among flow and diffusion models}

In the last two chapters we covered flow-matching, score-based and diffusion models. These were motivated from the common perspective of transforming a source of Gaussian noise into the target data distribution, but were presented as different generative procedures; however, we will see now that they are closely related. Based on an inspection of their corresponding loss functions, we will identify their similarities and the understand how their differences arise from the way their dynamics are parametrised.

\paragraph{Flow matching and score-based models.}

FM learns a deterministic velocity field $\v_\theta(\x_t,t)$ that approximates the true velocity field $\v^\star_t(\x_t)$ of a probability path $\{p_t\}_{t\in[0,1]}$. The objective is
\[
\mathcal{L}_{\text{FM}}
=
\mathbb{E}_{t,\x_t}
\left[
\|\v_\theta(\x_t,t) - \v^\star_t(\x_t)\|^2
\right].
\]
Then, the flow is given by the (learned) ODE
\[
\frac{d\x_t}{dt} = \v_\theta(\x_t,t),
\]
which transports samples from the source distribution to the data distribution.

Score-based models (SBM) instead learn the score function of the intermediate noise-corrupted densities,
\[
s_\theta(\x_t,t) \approx \nabla_{\x_t}\log p_t(\x_t),
\]
using a denoising score matching objective of the form
\[
\mathcal{L}_{\text{SBM}}
=
\mathbb{E}_{t,\x_t}
\left[
\|s_\theta(\x_t,t) - \nabla_{\x_t}\log p_t(\x_t)\|^2
\right].
\]

The connection between SBM and FM can be understood through the continuous-time formulation of score-based models introduced by \cite{song2021scorebased}. Here, the forward noising process is described by a stochastic differential equation
\[
d\x = f(\x,t)\,dt + g(t)\,d\w_t,
\]
where $f(\cdot)$ and $g(\cdot)$ are given by the noise perturbation schedule. 

As shown by \cite{song2021scorebased}, there exists a deterministic probability flow ODE associated with this diffusion process that induces the same marginal distributions $\{p_t\}_{t\in[0,1]}$. The velocity field of this ODE is given by
\[
\v_t^\star(\x_t)
=
f(\x_t,t)
-
\frac{1}{2} g(t)^2 \nabla_{\x_t}\log p_t(\x_t).
\]

Therefore, once the score function $\nabla_{\x_t}\log p_t(\x_t)$ is trained, the velocity field can be reconstructed. This means that SBMs implicitly learn the same transport vector field that FM methods parameterise directly.

\paragraph{Score-based models and DDPM.}

Denoising diffusion probabilistic models introduce the discrete-time forward noising process
\[
\x_t = \sqrt{\bar{\alpha}_t}\,\x_0 + \sqrt{1-\bar{\alpha}_t}\,\varepsilon,
\qquad \varepsilon \sim \mathcal{N}(0,I),
\]
and train a neural network to predict the noise $\varepsilon$ via
\[
\mathcal{L}_{\text{DDPM}}
=
\mathbb{E}_{t,\x_0,\varepsilon}
\left[
\|\varepsilon - \varepsilon_\theta(\x_t,t)\|^2
\right].
\]


Note that, since  the score of the Gaussian conditional $p(\x_t|\x_0)$ satisfies
\[
\nabla_{\x_t}\log p(\x_t|\x_0)
=
-\frac{1}{1-\bar{\alpha}_t}
(\x_t - \sqrt{\bar{\alpha}_t}\x_0),
\]
The SBM objective is equivalent to that of DDPM score matching up to a known scaling factor. Therefore, DDPM training can be interpreted as learning the time-dependent score $s_\theta(\x_t,t)$.

\paragraph{Flow matching and DDPM.}

Lastly, combining the previous equivalences implies that DDPM and FM also become equivalent.

From the perspective of the loss functions, though the methods adopt different representations, they are equivalent. FM regresses the velocity field $v^\star(\x_t,t)$, SBM regresses the score $\nabla\log p_t(\x_t)$, and DDPM regresses the noise variable $\varepsilon$. However, these quantities are linearly related given the shared assumption of the diffusion process. This means that  the three objectives are different parameterisations of the same concept, which is to learn a vector field that transports a Gaussian distribution to the data distribution over time.



%%%%%%%%%%%%%%%%%%%%%%%%%%%%%%%
%%%%%%%%%EXERCISES%%%%%%%%%%%%%
%%%%%%%%%%%%%%%%%%%%%%%%%%%%%%%


\clearpage
\newpage
\section{Exercises - Diffusion models}

\begin{exercise}[Score function properties]
Let $p(\x)$ be a smooth probability density on $\R^D$ with score $s(\x)=\nabla_\x \log p(\x)$.  
Show that $\mathbb{E}_{p(\x)}[s(\x)] = 0$ assuming that $p(\x)$ decays sufficiently fast at infinity, and compute $s(\x)$ explicitly when $p(\x)=\cN(\x;\mu,\Sigma)$, where $\mu\in\R^D$ and $\Sigma\in\R^{D\times D}$ is a symmetric positive definite matrix.
\end{exercise}

\begin{solution}[Score function properties]
Recall that
\[
s(\x)=\nabla_\x \log p(\x)=\frac{\nabla_\x p(\x)}{p(\x)},
\]
whenever $p(\x)>0$. Therefore,
\(
p(\x)s(\x)=\nabla_\x p(\x).
\)
Integrating with respect to $\x$ gives
\[
\mathbb{E}_{p(\x)}[s(\x)]
=
\int_{\R^D} s(\x)\,p(\x)\,\d\x
=
\int_{\R^D} \nabla_\x p(\x)\,\d\x.
\]
For each $j=1,\dots,D$, the $j$-th component of the rhs in the above expression is 
\(
\int_{\R^D} \partial_{x_j}p(\x)\,\d\x.
\)
Assuming $p(\x)$ decays at infinity, i.e.,  $p(\x)\to 0$ sufficiently fast as $\|\x\|\to\infty$, the fundamental theorem of calculus implies 
\[
\int_{\R^D} \partial_{x_j}p(\x)\,\d\x = 0.
\]
Hence every component vanishes, and thus
\(
\mathbb{E}_{p(\x)}[s(\x)] = 0.
\)

The log-density of
\(
p(\x)=\cN(\x;\mu,\Sigma),
\)
is 
\[
\log p(\x)
=
-\frac{D}{2}\log(2\pi)
-\frac{1}{2}\log\det(\Sigma)
-\frac{1}{2}(\x-\mu)^\top \Sigma^{-1}(\x-\mu).
\]
The first two terms are constant with respect to $\x$, so differentiating yields
\[
s(\x)=\nabla_\x \log p(\x)
=
-\frac{1}{2}\nabla_\x\left((\x-\mu)^\top \Sigma^{-1}(\x-\mu)\right).
\]
Since $\Sigma^{-1}$ is symmetric, we have that the score of the multivariate Gaussian distribution is 
\[
s(\x)= -\Sigma^{-1}(\x-\mu).
\]
Notice that, in this case in particular, we have 
\[
\mathbb{E}_{p(\x)}[s(\x)]
=
-\Sigma^{-1}\bigl(\mathbb{E}[\x]-\mu\bigr)
=0,
\]
which is consistent with the general result.
\end{solution}


\begin{exercise}[Langevin dynamics interpretation]
Let $p(\x)$ be a target density with score function $s(\x)=\nabla_\x \log p(\x)$.  
Consider the update
\[
\x_{k+1} = \x_k + \frac{\epsilon}{2}s(\x_k) + \sqrt{\epsilon}\,\eta_k,
\]
where $\epsilon>0$ is a step size and $\eta_k\sim\cN(0,I)$ is a standard Gaussian random vector.

Show that this corresponds to a stochastic gradient ascent step on $\log p(\x)$, and explain the role of the Gaussian noise term.
\end{exercise}

\begin{solution}[Langevin dynamics interpretation]
Since $s(\x)=\nabla_\x \log p(\x)$, the update can be written as
\[
\x_{k+1}
=
\x_k + \frac{\epsilon}{2}\nabla_\x \log p(\x_k) + \sqrt{\epsilon}\,\eta_k.
\]
Ignoring the noise term, this is a gradient ascent step on $\log p(\x)$ with step size $\epsilon/2$, so the deterministic part moves $\x_k$ toward regions of higher (log) probability density.

The parameter $\epsilon>0$ controls the magnitude of both the gradient step and the noise. Smaller values lead to more accurate but slower updates, while larger values allow for covering larger areas in the state space but may introduce instabilities.

The Gaussian noise term $\sqrt{\epsilon}\,\eta_k$, with $\eta_k\sim\cN(0,I)$, injects randomness into the dynamics, thus promoting exploration of the state space, preventing the chain locking on the local maxima.
\end{solution}


\begin{exercise}[Failure of score on manifolds]
Assume the data distribution is supported on a lower-dimensional manifold in $\R^D$.  
Explain why $p(\x)$ may not admit a density with respect to the Lebesgue measure, and why the score function becomes ill-defined.
\end{exercise}

\begin{solution}[Failure of score on manifolds]
If the data distribution is supported on a lower-dimensional manifold embedded in $\R^D$, then all its probability mass is concentrated on a set of Lebesgue measure zero in the ambient space. Therefore, the distribution is not absolutely continuous with respect to the Lebesgue measure, and so it does not admit a density $p(\x)$ on $\R^D$.

Since the score function is defined by
\[
s(\x)=\nabla_\x \log p(\x),
\]
it requires the existence of a differentiable density in the ambient space. When $p(\x)$ does not exist as a Lebesgue density, the expression $\log p(\x)$ is not well defined, and hence neither is its gradient. Therefore, the score function becomes ill-defined.
\end{solution}



\begin{exercise}[Score matching objective]
Starting from
\[
\cL_{\text{SM}}(\theta) = \frac{1}{2}\mathbb{E}_{p_{\text{data}}(\x)}\|s_\theta(\x) - \nabla_\x \log p_{\text{data}}(\x)\|^2,
\]
derive the equivalent expression
\[
\mathbb{E}_{p_{\text{data}}(\x)}\left[\tr(\nabla_\x s_\theta(\x)) + \frac{1}{2}\|s_\theta(\x)\|^2\right] + \text{constant}.
\]
\end{exercise}

\begin{solution}[Score matching objective]
Expanding the square gives
\begin{align*}
\cL_{\text{SM}}(\theta)
&=
\frac12 \int p_{\text{data}}(\x)
\left(
\|s_\theta(\x)\|^2
-2 s_\theta(\x)^\top \nabla_\x \log p_{\text{data}}(\x)
+\|\nabla_\x \log p_{\text{data}}(\x)\|^2
\right)\d\x \\
&=
\frac12 \int p_{\text{data}}(\x)\|s_\theta(\x)\|^2\,\d\x
-\int s_\theta(\x)^\top \nabla_\x p_{\text{data}}(\x)\,\d\x
+\text{constant},
\end{align*}
since
\[
p_{\text{data}}(\x)\nabla_\x \log p_{\text{data}}(\x)=\nabla_\x p_{\text{data}}(\x),
\]
and the last term does not depend on $\theta$.

Now write
\[
-\int s_\theta(\x)^\top \nabla_\x p_{\text{data}}(\x)\,\d\x
=
-\sum_{j=1}^D \int s_{\theta,j}(\x)\,\partial_{x_j}p_{\text{data}}(\x)\,\d\x.
\]
For each component, and assuming suitable decay at infinity, integration by parts gives
\[
-\int s_{\theta,j}(\x)\,\partial_{x_j}p_{\text{data}}(\x)\,\d\x
=
\int p_{\text{data}}(\x)\,\partial_{x_j}s_{\theta,j}(\x)\,\d\x.
\]
Summing over $j$ yields
\[
-\int s_\theta(\x)^\top \nabla_\x p_{\text{data}}(\x)\,\d\x
=
\int p_{\text{data}}(\x)\,\tr(\nabla_\x s_\theta(\x))\,\d\x.
\]

Therefore,
\[
\cL_{\text{SM}}(\theta)
=
\mathbb{E}_{p_{\text{data}}(\x)}
\left[
\tr(\nabla_\x s_\theta(\x))+\frac12\|s_\theta(\x)\|^2
\right]
+\text{constant}.
\]
\end{solution}



\begin{exercise}[Denoising score matching]
Let $q_\sigma(\tilde{\x}\mid\x)=\cN(\tilde{\x};\x,\sigma^2 I)$.  
Compute $\nabla_{\tilde{\x}}\log q_\sigma(\tilde{\x}\mid\x)$ and show that the DSM objective corresponds to a regression target $-(\tilde{\x}-\x)/\sigma^2$.
\end{exercise}

\begin{solution}[Denoising score matching]
Since
\[
q_\sigma(\tilde{\x}\mid\x)
=
\frac{1}{(2\pi\sigma^2)^{D/2}}
\exp\left(
-\frac{1}{2\sigma^2}\|\tilde{\x}-\x\|_2^2
\right),
\]
the score is given by 
\[
\nabla_{\tilde{\x}}\log q_\sigma(\tilde{\x}\mid\x)
=
-\frac{1}{\sigma^2}(\tilde{\x}-\x).
\]

Therefore, the denoising score matching objective
\[
\cL_{\mathrm{DSM}}(\theta)
=
\frac12
\mathbb{E}_{p_{\mathrm{data}}(\x)\,q_\sigma(\tilde{\x}\mid\x)}
\left\|
s_\theta(\tilde{\x})
-
\nabla_{\tilde{\x}}\log q_\sigma(\tilde{\x}\mid\x)
\right\|_2^2
\]
becomes
\[
\cL_{\mathrm{DSM}}(\theta)
=
\frac12
\mathbb{E}_{p_{\mathrm{data}}(\x)\,q_\sigma(\tilde{\x}\mid\x)}
\left\|
s_\theta(\tilde{\x})
+
\frac{\tilde{\x}-\x}{\sigma^2}
\right\|_2^2.
\]
This means that training $s_\theta$ is equivalent to solving a regression problem with the target
\(
-(\tilde{\x}-\x)/\sigma^2.
\)
\end{solution}


\begin{exercise}[Noise-conditional score network]
Let $\{\sigma_i\}_{i=1}^L$ be a collection of noise levels, and for each $\sigma_i$ define the denoising score matching loss
\[
l(\theta;\sigma_i)
=
\frac12
\mathbb{E}_{\x\sim p_{\mathrm{data}}(\x),\,\tilde{\x}\sim \cN(\x,\sigma_i^2 I)}
\left\|
s_\theta(\tilde{\x},\sigma_i)+\frac{\tilde{\x}-\x}{\sigma_i^2}
\right\|^2.
\]
The full training objective of the noise-conditional score network is
\[
\cL_{\mathrm{NCSN}}(\theta)
=
\frac1L\sum_{i=1}^L \lambda(\sigma_i)\,l(\theta;\sigma_i).
\]
Explain why choosing $\lambda(\sigma_i)=\sigma_i^2$ helps balance the contributions of different noise levels in the objective. Then, describe how one generates training pairs $(\tilde{\x},\sigma_i)$ in practice, without an explicit construction of $L$ separate noisy datasets.
\end{exercise}

\begin{solution}[Noise-conditional score network]
Recall from the previous exercise that, for a fixed noise level $\sigma_i$, the target in the loss is
\(
-(\tilde{\x}-\x)/\sigma_i^2.
\)
Since $\tilde{\x}=\x+\sigma_i \z$ with $\z\sim\cN(0,I)$, we have
\(
\tilde{\x}-\x=\sigma_i \z,
\)
and therefore the target becomes
\[
-\frac{\tilde{\x}-\x}{\sigma_i^2}=-\frac{\z}{\sigma_i}.
\]
Hence, its magnitude scales like $1/\sigma_i$, so the squared norm in the loss scales like $1/\sigma_i^2$. This means that, without weighting, the losses corresponding to small noise levels would tend to dominate the objective. Therefore, choosing
\[
\lambda(\sigma_i)=\sigma_i^2
\]
compensates for this different magnitudes, levelling the contributions from different noise levels.

In practice, one does not construct $L$ separate noisy datasets. Instead, during each training iteration, one first samples a clean datapoint
\(
\x\sim p_{\mathrm{data}}(\x), 
\)
and then (uniformly) samples a noise level $\sigma_i$ from the collection $\{\sigma_i\}_{i=1}^L$, to finally generate a noisy sample as
\[
\tilde{\x}=\x+\sigma_i \z,
\qquad
\z\sim\cN(0,I).
\]
The pair $(\tilde{\x},\sigma_i)$ is then used in the loss above. Using minibatches, we obtain a Monte Carlo approximation of the objective, without the need to store noisy copies of the dataset for every noise level.
\end{solution}

\begin{exercise}[Forward diffusion process]
Consider
\[
q(\x_t \mid \x_{t-1}) = \cN(\x_t;\sqrt{1-\beta_t}\x_{t-1},\beta_t I).
\]
Show that
\[
q(\x_t\mid\x_0)=\cN(\x_t;\sqrt{\bar\alpha_t}\x_0,(1-\bar\alpha_t)I),
\]
and derive $\bar\alpha_t=\prod_{s=1}^t(1-\beta_s)$.
\end{exercise}

\begin{solution}[Forward diffusion process]
Denoting $\alpha_t=1-\beta_t$, the forward transition can be expressed as 
\[
\x_t=\sqrt{\alpha_t}\x_{t-1}+\sqrt{1-\alpha_t}\,\epsilon_t,
\qquad \epsilon_t\sim\cN(0,I).
\]
Iterating this recursion gives
\[
\x_t
=
\sqrt{\alpha_t\alpha_{t-1}\cdots \alpha_1}\,\x_0
+
\text{Gaussian noise}.
\]
Hence, defining
\[
\bar\alpha_t=\prod_{s=1}^t \alpha_s=\prod_{s=1}^t(1-\beta_s),
\]
the mean of $\x_t\mid \x_0$ is
\[
\mathbb{E}[\x_t\mid \x_0]=\sqrt{\bar\alpha_t}\,\x_0.
\]

Since each step adds independent Gaussian noise, $\x_t\mid \x_0$ is Gaussian. Therefore, applying the covariance operator to the forward transition equation gives a recursion for covariance of $\x_t\mid \x_0$, denoted $\Sigma_t$, as
\[
\Sigma_t=\alpha_t \Sigma_{t-1}+(1-\alpha_t)I,
\qquad \Sigma_0=0,
\]
whose solution is
\[
\Sigma_t=(1-\bar\alpha_t)I.
\]
This gives the required result
\[
q(\x_t\mid\x_0)=\cN(\x_t;\sqrt{\bar\alpha_t}\x_0,(1-\bar\alpha_t)I).
\]

\end{solution}

\begin{exercise}[Noise prediction parametrisation]
In DDPM, the reverse process is parametrised as
\[
p_\theta(\x_{t-1}\mid \x_t)=\cN(\x_{t-1};\mu_\theta(\x_t,t),\sigma_t^2 I).
\]
Let $\epsilon_\theta(\x_t,t)$ be a neural network estimating the noise in $\x_t$, and consider
\[
\mu_\theta(\x_t,t) = \frac{1}{\sqrt{\alpha_t}}\left(\x_t - \frac{\beta_t}{\sqrt{1-\bar\alpha_t}}\epsilon_\theta(\x_t,t)\right).
\]
Show that, under the matching loss
\[
\|\tilde\mu_t(\x_t,\x_0)-\mu_\theta(\x_t,t)\|^2,
\]
the choice of  $\mu_\theta$ above leads to a loss of the form
\[
\|\epsilon - \epsilon_\theta(\x_t,t)\|^2.
\]
\end{exercise}

\begin{solution}[Noise prediction parametrisation]
Recall that
\[
\x_t=\sqrt{\bar\alpha_t}\x_0+\sqrt{1-\bar\alpha_t}\,\epsilon,
\qquad \epsilon\sim\cN(0,I),
\]
and that the true posterior mean of $q(\x_{t-1}\mid \x_t,\x_0)$ can be written as
\[
\tilde\mu_t(\x_t,\x_0)
=
\frac{1}{\sqrt{\alpha_t}}
\left(
\x_t-\frac{\beta_t}{\sqrt{1-\bar\alpha_t}}\epsilon
\right).
\]
Using the given parametrisation, we obtain
\[
\|\tilde\mu_t(\x_t,\x_0)-\mu_\theta(\x_t,t)\|^2
=
\frac{\beta_t^2}{\alpha_t(1-\bar\alpha_t)}
\|\epsilon-\epsilon_\theta(\x_t,t)\|^2.
\]
Thus, up to a known scalar weight depending only on $t$, the loss is of the form
\[
\|\epsilon-\epsilon_\theta(\x_t,t)\|^2.
\]
\end{solution}


